%% ma: L12-B02-0008
%% lop: 12
%% chuong: 1
%% bai: 2
%% dang: 12.2.1
%% loai: DS
%% muc: [NB, TH, TH, VD]
%% dap_an: "ĐSĐĐ"
%% nguon: "(NBV)-ĐỀ SỐ 10-MỖI NGÀY 1 ĐỀ THI 2025.pdf (D10, MỖI NGÀY 1 ĐỀ - Đề số 10), Phần II Câu 1"
%% kien_thuc: "Bài 2 (giá trị nhỏ nhất trên một đoạn), Bài 1 (cực trị), đạo hàm hàm lượng giác (lớp 11), phương trình lượng giác cơ bản (lớp 11)"
%% trang_thai: chinh-thuc
%% ghi_chu: "Cùng cấu trúc bốn ý với L12-B02-0001 (hàm lượng giác trên một đoạn: giá trị, đạo hàm, nghiệm f'(x) = 0, giá trị lớn nhất hoặc nhỏ nhất) nên nhập vào dạng 12.2.1 đã duyệt; cùng hàm số với L11-B32-0002 nhưng hỏi thêm giá trị nhỏ nhất. Đáp án gốc Đ S Đ Đ đúng (ý b sai vì f'(x) = 2cos 2x - 1)."
\begin{ex}
Cho hàm số $f(x)=\sin2x-x$. Xét tính đúng sai của các khẳng định sau.
\choiceTF
{\True $f\left(-\dfrac\pi2\right)=\dfrac\pi2$ và $f\left(\dfrac\pi2\right)=-\dfrac\pi2$.}
{Đạo hàm của hàm số đã cho là $f'(x)=\cos2x-1$.}
{\True Nghiệm của phương trình $f'(x)=0$ trên đoạn $\left[-\dfrac\pi2;\dfrac\pi2\right]$ là $-\dfrac\pi6$ hoặc $\dfrac\pi6$.}
{\True Giá trị nhỏ nhất của $f(x)$ trên đoạn $\left[-\dfrac\pi2;\dfrac\pi2\right]$ là $-\dfrac\pi2$.}
\loigiai{a) $f\left(-\dfrac\pi2\right)=\sin(-\pi)+\dfrac\pi2=\dfrac\pi2$; $f\left(\dfrac\pi2\right)=\sin\pi-\dfrac\pi2=-\dfrac\pi2$. Đúng.

b) $f'(x)=2\cos2x-1$. Sai.

c) $f'(x)=0\Leftrightarrow\cos2x=\dfrac12\Leftrightarrow x=\pm\dfrac\pi6+k\pi$ $(k\in\mathbb Z)$. Trên $\left[-\dfrac\pi2;\dfrac\pi2\right]$ có $x=-\dfrac\pi6$ và $x=\dfrac\pi6$. Đúng.

d) $f\left(\pm\dfrac\pi6\right)=\pm\left(\dfrac{\sqrt3}2-\dfrac\pi6\right)\approx\pm0{,}34$, $f\left(-\dfrac\pi2\right)=\dfrac\pi2$, $f\left(\dfrac\pi2\right)=-\dfrac\pi2\approx-1{,}57$. Vậy giá trị nhỏ nhất là $-\dfrac\pi2$. Đúng.}
\end{ex}
